% ═════════════════════════════════════════════════════════════════════════════
\section{Applications to Weak Memory Models}
\label{sec:wmm}
% ═════════════════════════════════════════════════════════════════════════════

Modern multiprocessor architectures do not guarantee Sequential Consistency
(SC)~\cite{Lam79}: for performance reasons, processors reorder memory
operations in ways that are invisible to a single thread but observable
between concurrently executing threads. The resulting \emph{weak} (or
\emph{relaxed}) \emph{memory models} (WMMs) include Total Store Order (TSO),
which underlies x86 and SPARC processors~\cite{OSS09,SSONM10}, and Partial
Store Order (PSO), among others~\cite{AMT14}. Verifying concurrent programs
under such models is substantially harder than under SC, because each thread
is equipped with a buffer that decouples writes from their visibility to other
threads, making the state space infinite even for finite-state programs.
Parameterized verification under WMMs is harder still: one must reason
simultaneously about an unbounded number of threads, each carrying an
unbounded buffer.

Parameterized questions arise naturally in this setting. Synchronization
primitives, such as locks, barriers, and mutual exclusion protocols, are
written to be used by any number of threads, and the same code is executed by
all of them. A correct implementation must work for every number of threads,
and on every architecture on which it is deployed, which means under the
memory model of that architecture.

The connection to the techniques covered in earlier sections is direct.
Backward reachability (Section~\ref{sec:wqo}) applies when the system's
transition relation is monotone with respect to a suitable
well-quasi-ordering on configurations, and monotonic abstraction
(Section~\ref{sec:rmc}) enforces monotonicity by over-approximation when it
does not hold. The key contribution of the weak-memory line of work is to identify the right
semantics and the right ordering to make this apply, even though store and
load buffers introduce data that naively breaks monotonicity.

\subsection{The TSO Model and the Dual Semantics}

Under the classical TSO semantics, each thread is equipped with an unbounded
perfect FIFO \emph{store buffer}. When a thread writes to a variable, the
write is appended to its store buffer rather than committed to memory
immediately; values propagate to shared memory non-deterministically in FIFO
order. A thread reading a variable first checks its own buffer for the most
recent pending write, and falls back to main memory if none is found. Fence
instructions force the buffer to drain before execution continues.

\begin{figure}[t]
\centering
\begin{tabular}{l@{\hspace{3em}}l}
\toprule
Thread 1 & Thread 2 \\
\midrule
$x := 1$ & $y := 1$ \\
$r_1 := y$ & $r_2 := x$ \\
\bottomrule
\end{tabular}
\caption{The store buffering program. The variables $x$ and~$y$ are initially
$0$. Under SC, at least one of the two reads returns~$1$. Under TSO, both reads
may return~$0$.}
\label{fig:sb}
\end{figure}

The program of Figure~\ref{fig:sb} illustrates the difference from SC. Under
SC, one of the two writes is executed first, so at least one of the reads
returns the value~$1$. Under TSO, thread~1 can place its write to~$x$ in its
store buffer, and thread~2 can do the same with its write to~$y$. Each thread
then reads the other variable from memory, where it still has the value~$0$,
so the read of~$y$ by thread~1 returns~$0$, and the read of~$x$ by thread~2
returns~$0$. Only afterwards are the two writes propagated to memory. This
behavior breaks mutual exclusion protocols such as Dekker's and Peterson's
algorithms, which rely on the fact that a thread that has announced its
intention by a write will be seen by the others. Fences are needed to restore
correctness, and verification under TSO is used, among other things, to check
that enough fences have been inserted. We say that a program is
\emph{robust} if all its behaviors under TSO are also possible under SC.
Checking robustness is a related verification problem~\cite{BDM13}.

Decidability of safety verification under TSO for a fixed number of threads
was established by encoding store buffers as lossy channels and applying the
theory of well-structured transition systems (WSTS)~\cite{ABBM10}. However,
these encodings are complex: each channel message must carry a full memory
snapshot, and the reduction involves non-deterministic guessing that leads to
significant state-space explosion in practice.

An alternative approach introduces a \emph{dual} semantics for
TSO~\cite{AABN16,AABN18}. Rather than store buffers, the dual model
equips each thread with a \emph{load buffer}: values propagate
non-deterministically from shared memory into load buffers, and reads are
served from the head of the load buffer rather than from memory directly.
Writes, in contrast, are committed to memory immediately and tagged in the
buffer as \emph{own} messages. In the program of Figure~\ref{fig:sb}, the
initial values of $x$ and~$y$ can propagate into the load buffers of
threads~2 and~1, respectively, before any write takes place. Both writes are
then committed to memory, and each thread reads the old value~$0$ from its
load buffer. Thus, instead of delaying the writes, the dual semantics delays
the reads. The dual semantics is provably equivalent to the classical one, in
the sense that the two models reach the same sets of local states. It offers
significant practical advantages. Because load buffers behave as lossy channels in a
straightforward sense, the WSTS framework applies with far less overhead than
in the store-buffer encoding, yielding a simpler decidability proof and a more
efficient verification algorithm.

\subsection{Parameterized Verification under TSO}

The dual-semantics approach extends naturally to parameterized verification.
In the parameterized setting, the number of threads is a free parameter, and
one asks whether a bad local state is reachable for \emph{any} number of
threads. Parameterized configurations are ordered by a well-quasi-ordering
that combines the subword ordering on load buffer contents with a matching on
threads, in which each thread of the smaller configuration is mapped to a
distinct thread of the larger one; the transition system is shown to be
monotone with respect to this ordering; and backward reachability over
upward-closed sets of configurations terminates by Higman's lemma. This yields
the first decidability result for parameterized verification of concurrent
programs under a weak memory model~\cite{AABN16,AABN18}. Moreover, experiments
with a prototype tool show that parameterized verification is often
\emph{more} efficient than bounding the number of threads to a concrete
value: the sets of minimal elements produced by the backward analysis are
typically much smaller than the configuration sets generated by
bounded-instance analyses. Already for a fixed number of threads, the
load-buffer semantics gives large gains in efficiency and scalability compared
with verification tools based on the store-buffer semantics, on standard mutual
exclusion and memory-model benchmarks~\cite{AABN18}.

The observation that load buffers, unlike store buffers, are compatible with
monotonicity is the central point. Adding a message to a store buffer can
change the value that its thread reads, since the thread reads its own most
recent write, and can therefore disable a transition. Adding a message to a
load buffer only adds the possibility of reading an old value, which the
thread is free to skip. Larger configurations can thus simulate smaller ones,
as required by the theory of Section~\ref{sec:wqo}.

\paragraph{Complexity of parameterized TSO verification.}
While the dual-semantics approach establishes decidability, it leaves open
the precise computational complexity of parameterized TSO verification. This
question was settled in~\cite{AAR20}, where it is shown that the reachability
problem for a parameterized system of identical processes under TSO is
\textsc{PSpace}-complete.

The \textsc{PSpace} upper bound is obtained by combining the load-buffer
(dual) semantics with a careful observation about what information needs to be
tracked: it suffices to maintain, for each process, the set of write
operations it has issued together with the current value it observes for each
variable, rather than the full buffer contents. This yields a polynomial-space
abstraction that is both sound and complete for reachability. The matching
lower bound is established by a reduction from the reachability problem for
1-safe Petri nets, and the reduction relies on the ability of
processes to read the \emph{initial values} of variables from memory. This
dependence on initial values is not merely an artifact of the proof: when
processes are forbidden from reading initial memory values, the reachability
problem drops to polynomial time, since the relevant information can be
captured by a polynomial-sized abstraction tracking only the set of writes
issued across all processes. A further variant is studied in which processes
may perform \emph{atomic read-modify-write} (RMW) operations; it is shown
that adding RMW makes the problem non-primitive recursive,
matching the known hardness of reachability for general lossy channel systems.
Together, these results give a precise and tight complexity picture for
parameterized TSO verification across different fragments of the model.
Parameterized verification has also been studied for other memory models; for
instance, it is \textsc{PSpace}-complete under the release-acquire semantics
of C11~\cite{KGMC22}.

\subsection{The Intel x86 (eTSO) Model}

Real Intel x86 processors implement a memory model that goes beyond TSO. The
formal model eTSO~\cite{RMV22} accounts for multiple memory types, namely
write-back (wb), write-combining (wc), and strong-uncacheable, as well as
non-temporal writes (ntw). Write-back writes behave as in TSO, while
non-temporal writes on different variables may overtake one another, similarly
to PSO. This makes eTSO strictly more relaxed than TSO, and the interaction
between wb and ntw operations introduces subtle orderings that do not arise in
either TSO or PSO alone. Reachability under eTSO for a fixed number of
threads is in general undecidable, but decidability can be recovered by
bounding the number of \emph{alternations} between wb and ntw write stages
per thread~\cite{AABKS24}. A run is $k$-alternation bounded if the writes of
each thread can be partitioned into at most $k$ blocks, alternating between
blocks containing only wb writes and blocks containing only ntw writes. Even
under this restriction, executions may involve unbounded sequences of
operations within each block, and complex shuffles of writes across blocks are
permitted.

Parameterized verification of this $k$-alternation bounded problem was
established in~\cite{AABKS26}. The proof introduces a new transition system,
called the \emph{NTWLoad} system, which combines the load-buffer technique
from the dual TSO semantics (to handle wb writes while preserving monotonicity)
with per-variable store buffers for ntw writes. Each eTSO thread is simulated
by $2k$ concurrent \emph{stage threads}, alternating between wb and ntw
stages, with a finite interface that enforces the required causal dependencies
across stages.

A central technical difficulty is, as explained above, that TSO store buffers
are not monotone with respect to the subword ordering. This difficulty is circumvented by switching wb writes to the
load-buffer representation, in which monotonicity is restored, while ntw
writes retain their store-buffer treatment. The causal dependencies across
stages are maintained through an explicit bookkeeping mechanism. For each
stage, it tracks speculated write values, expected termination labels, and
obligations that must be discharged during execution.

This bookkeeping mechanism introduces a further complication: the resulting
ordering on configurations is not merely a well-quasi-ordering (WQO) but must
be a \emph{better quasi-ordering} (BQO)~\cite{Mar94}. BQOs are strictly
stronger than WQOs and enjoy additional closure properties, in particular
under the powerset construction, that WQOs do not. The ordering on NTWLoad
configurations is built by combining a subword ordering on buffer contents
with an equality constraint on stage information and a powerset-style matching
on thread sets; showing this composite ordering is a BQO requires appealing to
the closure of BQOs under products, word embeddings, and powerset operations.
The transition relation is then shown to be monotone with respect to this BQO,
placing the system within the WSTS framework and yielding decidability of
backward reachability. BQOs had earlier been used in verification to obtain
efficient symbolic representations for timed Petri nets~\cite{AN00,Abd10}. To
the best of our knowledge, the result of~\cite{AABKS26} is the first
application of BQO-based well-structured systems to prove decidability of a
program verification problem.
